\documentclass[11pt,a4paper]{article}
\usepackage[T1]{fontenc}
\usepackage[utf8]{inputenc}
\usepackage{amsmath,amssymb,amsthm}
\usepackage[margin=1in]{geometry}
\usepackage[colorlinks=true,linkcolor=blue,urlcolor=blue]{hyperref}
\theoremstyle{plain}
\newtheorem{theorem}{Theorem}
\newtheorem{lemma}[theorem]{Lemma}
\newtheorem{proposition}[theorem]{Proposition}
\newtheorem{corollary}[theorem]{Corollary}
\theoremstyle{definition}
\newtheorem{remark}[theorem]{Remark}

\title{The empirical recurrence for OEIS A249962, proved from an exact model}
\author{Adrian Perez Fontelles\\ \small Independent researcher}
\date{2 September 2026}
\begin{document}
\maketitle

\begin{abstract}
OEIS A249962 carries an empirical recurrence of order $8$. It is true, and it is
decidable rather than empirical. The entry's sequence is given exactly by an exact polynomial in $n$, from
which $a$ provably satisfies a MONIC linear recurrence of order $S=8$, derived below and not
assumed. The residual of the conjectured recurrence satisfies the same one, so whether it
annihilates $a$ is settled by a finite exact computation. No transfer matrix is involved and
the count is not a walk.
\end{abstract}

\noindent\small 2020 Mathematics Subject Classification. 05A15, 05A19, 11P21.\normalsize

\section{The conjecture}

OEIS A249962 is ``Number of length 2+5 0..n arrays with no six consecutive terms having the maximum of any two terms equal to the minimum of the remaining four terms.''. It has offset $1$ and begins
\[
15, 505, 5300, 31180, 129095, 422065, 1164800, 2830080,\ \dots
\]
The entry states, as a conjecture contributed by its author and never marked settled:
\begin{quote}
Empirical: a(n) = n\^{}7 + (13/6)*n\^{}6 + (13/3)*n\^{}5 + (65/12)*n\^{}4 + (3/2)*n\^{}3 - (1/12)*n\^{}2 + (2/3)*n.
\end{quote}
As of the ``Last modified'' line on the live entry (Aug 21 2017 08:22:40, revision 8)
this is still recorded as empirical, and nothing on the entry records it as proved.

\section{The count is a polynomial}

The array has a FIXED length $L$ and an alphabet $0..n$ that grows. The condition is decided by
the ORDER of the terms alone --- by which are equal and which is larger --- so an admissible
array is described completely by its weak ordering: the ordered set partition of the $L$
positions into $m$ blocks of equal value, blocks listed in increasing order. Every assignment
of $m$ distinct values from $\{0,\dots,n\}$ to those blocks, in increasing order, gives one
array, and there are $\binom{n+1}{m}$ of them. Hence
\[
a(n)\;=\;\sum_{m=1}^{L} N_m\binom{n+1}{m},
\]
with $N_m$ the number of admissible weak orderings into $m$ blocks, computed exactly.

\section{The bound}

That expression is a polynomial in $n$ of degree at most $L$, so $(z-1)^{L+1}$ annihilates
$a$ and $S=8$. Nothing is fitted and no threshold is guessed: the $N_m$ are counted, not
estimated.

\section{The proof}

\begin{lemma}\label{lem:crit}
Let $A(z)=z^S-\sum_{j<S}\alpha_jz^{j}$ be the monic annihilator derived above and
$q(z)=z^{8}-\sum_ic_iz^{8-i}$ the characteristic polynomial of the conjectured
recurrence. The residual $r(n)=a(n)-\sum_ic_ia(n-i)$ is a linear combination of shifts of $a$,
so $A$ annihilates $r$ as well. Since $A(0)\ne0$ the recurrence it gives runs in both
directions, and $S$ consecutive zeros of $r$ force $r$ to vanish at every later index.
\end{lemma}

\begin{proof}
$A$ annihilates $a$, hence every shift of $a$, hence any linear combination of shifts; that is
$r$. A monic recurrence with nonzero constant term determines each term from the $S$ before it
and each from the $S$ after it, so a run of $S$ zeros propagates.
\end{proof}

\begin{theorem}
\[
a(n)\;=\;8\,a(n-1) - 28\,a(n-2) + 56\,a(n-3) - 70\,a(n-4) + 56\,a(n-5) - 28\,a(n-6) + 8\,a(n-7) - a(n-8)
\]
for every $n>0$.
\end{theorem}

\begin{proof}
The terms $a(0),\dots$ were computed exactly in integer arithmetic from the model, extended by
$A$ where the model itself was not evaluated, and $r(n)$ formed for each index. Those integers
vanish from $n=0+1$ onwards, and the run exceeds $S+8$, which by
Lemma~\ref{lem:crit} settles every larger index at once.
\end{proof}

\section{Verification}

Three checks, each independent of the algebra above.

First, the model reproduces all $23$ terms the entry publishes, exactly, in integer
arithmetic. This is what ties the model to the entry: it is built by reading the entry's
English, and a misreading gives different counts.

Second, the conjectured recurrence was evaluated directly on those published terms, with no
model involved, and holds wherever the proved range and the data overlap.

Third, the derived annihilator $A$ was itself tested on terms it had not been given: the model
was evaluated beyond the $S$ values that determine it and $A$ reproduced them. A bound that is
too small fails this test, which is why it is run.

\begin{thebibliography}{9}
\bibitem{oeis} The OEIS Foundation, \emph{The On-Line Encyclopedia of Integer
Sequences}, \url{https://oeis.org}, 2026. Sequence A249962.
\bibitem{beckrobins} M.~Beck and S.~Robins, \emph{Computing the Continuous Discretely},
2nd ed., Springer, 2015. (Ehrhart quasi-polynomials and their periods.)
\bibitem{stanley} R.~P.~Stanley, \emph{Enumerative Combinatorics}, Volume 1, 2nd ed.,
Cambridge University Press, 2011.
\bibitem{ds} F.~Diamond and J.~Shurman, \emph{A First Course in Modular Forms},
Springer GTM 228, 2005.
\end{thebibliography}

\end{document}
